%2multibyte Version: 5.50.0.2953 CodePage: 1252

\documentclass[notes=show]{beamer}
%%%%%%%%%%%%%%%%%%%%%%%%%%%%%%%%%%%%%%%%%%%%%%%%%%%%%%%%%%%%%%%%%%%%%%%%%%%%%%%%%%%%%%%%%%%%%%%%%%%%%%%%%%%%%%%%%%%%%%%%%%%%%%%%%%%%%%%%%%%%%%%%%%%%%%%%%%%%%%%%%%%%%%%%%%%%%%%%%%%%%%%%%%%%%%%%%%%%%%%%%%%%%%%%%%%%%%%%%%%%%%%%%%%%%%%%%%%%%%%%%%%%%%%%%%%%
\usepackage{amsfonts}
\usepackage{graphicx}
\usepackage{amssymb}
\usepackage{amsmath}
\usepackage{mathpazo}
\usepackage{hyperref}
\usepackage{multimedia}

\setcounter{MaxMatrixCols}{10}
%TCIDATA{OutputFilter=LATEX.DLL}
%TCIDATA{Version=5.50.0.2953}
%TCIDATA{Codepage=1252}
%TCIDATA{<META NAME="SaveForMode" CONTENT="2">}
%TCIDATA{BibliographyScheme=Manual}
%TCIDATA{Created=Wednesday, May 01, 2013 14:54:39}
%TCIDATA{LastRevised=Wednesday, November 20, 2019 14:27:54}
%TCIDATA{<META NAME="GraphicsSave" CONTENT="32">}
%TCIDATA{<META NAME="DocumentShell" CONTENT="Other Documents\SW\Slides - Beamer">}
%TCIDATA{Language=American English}
%TCIDATA{CSTFile=beamer.cst}
%TCIDATA{PageSetup=72,72,72,83,0}
%TCIDATA{Counters=arabic,1}
%TCIDATA{AllPages=
%H=36
%F=36
%}


\newenvironment{stepenumerate}{\begin{enumerate}[<+->]}{\end{enumerate}}
\newenvironment{stepitemize}{\begin{itemize}[<+->]}{\end{itemize} }
\newenvironment{stepenumeratewithalert}{\begin{enumerate}[<+-| alert@+>]}{\end{enumerate}}
\newenvironment{stepitemizewithalert}{\begin{itemize}[<+-| alert@+>]}{\end{itemize} }
\input{tcilatex}
\usetheme{JuanLesPins}
\usecolortheme{whale}

\begin{document}

\title{Econometrics: Lectures 7\\
Introduction to Maximum Likelihood Estimation}
\author{Katerina Petrova \ \ \ \ \ \ \ \ \ \ \ \ \ \ \ \ \ \ \ \ \ \ \ \ \ \
\ \ \ \ \ \ \ \ \ \ \ \ \ \ \ \ \ \ \ \ \ \ \ \ \ \ \ \ \ \ \ \ \ \ \ \ \ \
\ \ \ \ \ \ \ \ \ \ \ \ \ \ \ \ \ \ \ \ \ \ \ }
\date{}
\maketitle

\section{The Likelihood Principle}

%TCIMACRO{\TeXButton{BeginFrame}{\begin{frame}}}%
%BeginExpansion
\begin{frame}%
%EndExpansion

\QTR{frametitle}{The Method of Maximum Likelihood}

\begin{itemize}
\item Suppose that $X_{1},X_{2},...,X_{n}$ are random variables from a
discrete family of distributions $\left\{ f\left( X_{1},..,X_{n},\theta
\right) :\theta \in \Theta \right\} $.

\item The probability that $X_{1}=x_{1},..$, $X_{n}=x_{n}$ is the joint mass
function:%
\begin{equation*}
f_{X}\left( x_{1},...,x_{n},\theta \right) =P_{\theta }\left(
X_{1}=x_{1},...,X_{n}=x_{n}\right) .
\end{equation*}

\item We might ask what value of $\theta $ maximises the probability of
obtaining this observed particular sample $x_{1},..,x_{n}$, i.e., what value
of $\theta $ maximises%
\begin{equation*}
f_{X}\left( x_{1},...,x_{n},\theta \right) .
\end{equation*}

\item This maximising value would provide a good estimator of $\theta $,
because it would provide the largest probability of this particular sample
occurring.

\item Put differently, the maximising value of $\theta $ picks up the most
likely mass to have generated $x_{1},...,x_{n}$.
\end{itemize}

%TCIMACRO{\TeXButton{EndFrame}{\end{frame}}}%
%BeginExpansion
\end{frame}%
%EndExpansion

%TCIMACRO{\TeXButton{BeginFrame}{\begin{frame}}}%
%BeginExpansion
\begin{frame}%
%EndExpansion

\QTR{frametitle}{The Method of Maximum Likelihood}

\begin{itemize}
\item In general $X=\left( X_{1},...,X_{n}\right) ^{\prime }$ is a random
vector from a family of densities (or mass functions) $\left\{ f\left(
X,\theta \right) :\theta \in \Theta \right\} $.
\end{itemize}

\begin{definition}
The likelihood function, $L\left( \theta \right) $ is defined as the density
(or mass) of the data $X=\left( X_{1},...,X_{n}\right) ^{\prime }$, but
regarded as function of $\theta $ rather than $X:$%
\begin{equation*}
L\left( \theta \right) =f\left( X,\theta \right) .
\end{equation*}%
The maximum likelihood estimator (MLE) $\hat{\theta}_{n}$ of $\theta $ is
defined as the maximiser of the likelihood function $L\left( \theta \right) $
over the parameter space $\Theta $, i.e., 
\begin{equation*}
\hat{\theta}_{n}=\arg \max_{\theta \in \Theta }L\left( \theta \right) \text{%
\ \ \ or\ \ \ }L\left( \hat{\theta}_{n}\right) =\max_{\theta \in \Theta
}L\left( \theta \right) .
\end{equation*}
\end{definition}

%TCIMACRO{\TeXButton{EndFrame}{\end{frame}}}%
%BeginExpansion
\end{frame}%
%EndExpansion

%TCIMACRO{\TeXButton{BeginFrame}{\begin{frame}}}%
%BeginExpansion
\begin{frame}%
%EndExpansion

\QTR{frametitle}{The Method of Maximum Likelihood}

\begin{itemize}
\item We define the log-likelihood function%
\begin{equation*}
l\left( \theta \right) :=\log L\left( \theta \right)
\end{equation*}

\item The MLE of $\theta $ can EQUIVALENTLY be found as the maximiser of $%
l\left( \theta \right) :$%
\begin{equation*}
\hat{\theta}_{n}=\arg \max_{\theta \in \Theta }l\left( \theta \right) \text{%
\ \ \ or\ \ \ }l\left( \hat{\theta}_{n}\right) =\max_{\theta \in \Theta
}l\left( \theta \right) .
\end{equation*}

\item Simply an algebraic convenience to use $l\left( \theta \right) $.
\end{itemize}

%TCIMACRO{\TeXButton{EndFrame}{\end{frame}}}%
%BeginExpansion
\end{frame}%
%EndExpansion

%TCIMACRO{\TeXButton{BeginFrame}{\begin{frame}}}%
%BeginExpansion
\begin{frame}%
%EndExpansion

Let's proof that if $h\left( \cdot \right) $\ is a positive and
differentiable function, $h\left( x\right) $\ is maximised at the same point
as $\log h\left( x\right) $\emph{.}

\begin{proof}
Take $x^{\ast }$ to be the maximum of $\log h\left( x\right) .$ That means,
that $\frac{\partial \log h\left( x\right) }{\partial x}|_{x=x^{\ast }}=0$
since $x^{\ast }$ is a maximum and $h\left( \cdot \right) $ is
differentiable.%
\begin{equation*}
\Longrightarrow \frac{\partial \log h\left( x\right) }{\partial x}%
|_{x=x^{\ast }}=\frac{1}{h\left( x\right) }\frac{\partial h\left( x\right) }{%
\partial x}|_{x=x^{\ast }}=0
\end{equation*}%
by the chain rule for differentiation and since $h\left( x\right) $ is
positive, 
\begin{equation*}
\Longrightarrow \frac{\partial h\left( x\right) }{\partial x}|_{x=x^{\ast
}}=0
\end{equation*}%
Moreover, $\frac{\partial h\left( x\right) }{\partial x}$ and $\frac{%
\partial \log h\left( x\right) }{\partial x}$ have the same sign, so $%
h\left( \cdot \right) $ increases when $\log h\left( \cdot \right) $
increases, and vice versa so they have the same maxima.
\end{proof}

%TCIMACRO{\TeXButton{EndFrame}{\end{frame}}}%
%BeginExpansion
\end{frame}%
%EndExpansion

%TCIMACRO{\TeXButton{BeginFrame}{\begin{frame}}}%
%BeginExpansion
\begin{frame}%
%EndExpansion

\QTR{frametitle}{The Method of Maximum Likelihood}

\begin{lemma}
If the random variable $X$ has density $f\left( X,\theta _{0}\right) $ with
support $\chi $ and $\int_{\chi }f\left( y,\theta \right) dy=1$ for all $%
\theta \in \Theta $, then the function $g\left( \theta \right) =\mathbb{E}%
\log f\left( X,\theta \right) $ is maximised at $\theta =\theta _{0}$.
\end{lemma}

%TCIMACRO{\TeXButton{EndFrame}{\end{frame}}}%
%BeginExpansion
\end{frame}%
%EndExpansion

%TCIMACRO{\TeXButton{BeginFrame}{\begin{frame}}}%
%BeginExpansion
\begin{frame}%
%EndExpansion

\QTR{frametitle}{The Method of Maximum Likelihood}

\begin{proof}
By the Jensen inequality for concave functions, $\mathbb{E}\left( \log
Y\right) \leq \log \mathbb{E}\left( Y\right) $ for any random variable $Y$.
Now write%
\begin{eqnarray*}
&&\mathbb{E}\log f\left( X,\theta \right) -\mathbb{E}\log f\left( X,\theta
_{0}\right) \\
&=&\mathbb{E}\left[ \log \frac{f\left( X,\theta \right) }{f\left( X,\theta
_{0}\right) }\right] \leq \log \mathbb{E}\left[ \frac{f\left( X,\theta
\right) }{f\left( X,\theta _{0}\right) }\right] \;\;\;\;\left( \text{Jensen}%
\right) \\
&=&\log \int_{\chi }\frac{f\left( X,\theta \right) }{f\left( X,\theta
_{0}\right) }f\left( X,\theta _{0}\right) dX \\
&=&\log \int_{\chi }f\left( X,\theta \right) dX=\log 1=0
\end{eqnarray*}%
since $\theta \in \Theta ,$ so $\int_{\chi }f\left( X,\theta \right) dX=1$.
So, $\mathbb{E}\log f\left( X,\theta \right) \leq \mathbb{E}\log f\left(
X,\theta _{0}\right) $ for all $\theta \in \Theta .$
\end{proof}

%TCIMACRO{\TeXButton{EndFrame}{\end{frame}}}%
%BeginExpansion
\end{frame}%
%EndExpansion

%TCIMACRO{\TeXButton{BeginFrame}{\begin{frame}}}%
%BeginExpansion
\begin{frame}%
%EndExpansion

\QTR{frametitle}{The Cram\'{e}r Rao Lower Bound}

\begin{theorem}[Cram\'{e}r-Rao]
Under some regularity conditions, the covariance matrix of any unbiased
estimator $\hat{\theta}$ of $\theta _{0}$\ satisfies the following bound%
\begin{equation*}
V\left( \hat{\theta}\right) \geq \left[ \mathcal{I}\left( \theta _{0}\right) %
\right] ^{-1}
\end{equation*}%
in the sense that $V\left( \hat{\theta}\right) -\left[ \mathcal{I}\left(
\theta _{0}\right) \right] ^{-1}$ is a positive semi-definite matrix, where $%
\mathcal{I}\left( \theta \right) =-\mathbb{E}\left[ \frac{\partial ^{2}\log
L(\theta )}{\partial \theta \partial \theta ^{\prime }}\right] .$
\end{theorem}

\begin{itemize}
\item The theorem implies that the MSE of any unbiased estimator $\hat{\theta%
}$ of $\theta _{0}$ cannot be \textquotedblleft smaller\textquotedblright\
than $\left[ \mathcal{I}\left( \theta _{0}\right) \right] ^{-1}$, i.e. the
best possible unbiased estimator is one whose covariance matrix achieves the
Cram\'{e}r Rao lower bound. Such estimators are called efficient.
\end{itemize}

%TCIMACRO{\TeXButton{EndFrame}{\end{frame}}}%
%BeginExpansion
\end{frame}%
%EndExpansion

%TCIMACRO{\TeXButton{BeginFrame}{\begin{frame}}}%
%BeginExpansion
\begin{frame}%
%EndExpansion

\QTR{frametitle}{The Method of Maximum Likelihood}

\begin{definition}
An unbiased estimator $\hat{\theta}_{n}$ of $\theta _{0}$ is called \textbf{%
efficient} if 
\begin{equation*}
V\left( \hat{\theta}_{n}\right) =\left[ \mathcal{I}\left( \theta _{0}\right) %
\right] ^{-1}.
\end{equation*}
\end{definition}

\begin{itemize}
\item MLE estimators are consistent (under some regularity conditions) 
\begin{equation*}
\hat{\theta}_{n}\rightarrow _{p}\theta _{0}\text{ \ as \ }n\rightarrow
\infty .
\end{equation*}

\item And asymptotically normal:%
\begin{equation*}
\sqrt{n}(\hat{\theta}_{n}-\theta _{0})\!\approx \text{\ }\mathcal{N}\left(
0,n\left[ \mathcal{I}\left( \theta _{0}\right) \right] ^{-1}\right)
\end{equation*}

\item Finally, the MLE\ achieves asymptotically the CR Lower Bound, i.e.
it's asymptotically efficient.
\end{itemize}

%TCIMACRO{\TeXButton{EndFrame}{\end{frame}}}%
%BeginExpansion
\end{frame}%
%EndExpansion

%TCIMACRO{\TeXButton{BeginFrame}{\begin{frame}}}%
%BeginExpansion
\begin{frame}%
%EndExpansion

\QTR{frametitle}{The Method of Maximum Likelihood}

\begin{itemize}
\item An interesting result, that we will not prove is that:

\item If an MLE\ estimator is unbiased and efficient, then:

\begin{enumerate}
\item It is unique.

\item It necessarily belongs to the exponential family of distributions.
\end{enumerate}
\end{itemize}

\begin{definition}
The exponential family of distributions is a set of probability
distributions of the form $f(X,\theta )=h(X)g(\theta )\exp (\eta (\theta
)\cdot T(X))$
\end{definition}

\begin{itemize}
\item The family includes many of the most common distributions: Normal,
Chi-squared, Dirichlet, Bernoulli, Poisson, Gamma.

\item Since the exponential family is widely used, you can see why working
with the $log$ of the likelihood makes sense

\item Examples of exceptions are T-distribution, Uniform, most mixture
distributions.
\end{itemize}

%TCIMACRO{\TeXButton{EndFrame}{\end{frame}}}%
%BeginExpansion
\end{frame}%
%EndExpansion

%TCIMACRO{\TeXButton{BeginFrame}{\begin{frame}}}%
%BeginExpansion
\begin{frame}%
%EndExpansion

\QTR{frametitle}{The Method of Maximum Likelihood}

\begin{itemize}
\item How is this different from doing OLS?

\begin{enumerate}
\item OLS is a linear model, MLE\ can take any form and distribution.

\item In most applications, there is not even an analytic expression, so
numerical methods must be used.

\item OLS can be regarded as semi-parametric, the distribution of $%
\varepsilon $ is not specified.

\item MLE is parametric, we need to specify a distribution. This may be too
restrictive.

\item MLE, therefore, achieves the parametric Lower Bound, that is CRLB.

\item Under mis-specified distribution, one can still get consistency, i.e.
Quasi-MLE, but not efficiency.
\end{enumerate}
\end{itemize}

%TCIMACRO{\TeXButton{EndFrame}{\end{frame}}}%
%BeginExpansion
\end{frame}%
%EndExpansion

\section{MLE\ Examples}

%TCIMACRO{\TeXButton{BeginFrame}{\begin{frame}}}%
%BeginExpansion
\begin{frame}%
%EndExpansion

\QTR{frametitle}{The Method of Maximum Likelihood}

\begin{definition}[Independent sampling]
When $X=\left( X_{1},...,X_{n}\right) ^{\prime }$\ consists of independent
random variables $X_{1},...,X_{n}$,%
\begin{equation*}
f\left( X_{1},...,X_{n}\right) =f\left( X_{1}\right) ...f\left( X_{n}\right)
.
\end{equation*}%
Therefore,%
\begin{equation*}
L\left( \theta \right) =\prod\limits_{i=1}^{n}f\left( X_{i},\theta \right)
\quad \text{and}\quad l\left( \theta \right) =\sum_{i=1}^{n}\log f\left(
X_{i},\theta \right) .
\end{equation*}
\end{definition}

%TCIMACRO{\TeXButton{EndFrame}{\end{frame}}}%
%BeginExpansion
\end{frame}%
%EndExpansion

%TCIMACRO{\TeXButton{BeginFrame}{\begin{frame}}}%
%BeginExpansion
\begin{frame}%
%EndExpansion

\QTR{frametitle}{The Method of Maximum Likelihood}

\begin{example}[Normal with known variance]
Suppose $X_{1},...,X_{n}$ is an i.i.d. sample of a Normal Distribution $%
\mathcal{N}\left( \mu ,1\right) $ with density function: $f(X_{i};\mu )=%
\frac{1}{\sqrt{2\pi }}e^{-\frac{\left( X_{i}-\mu \right) ^{2}}{2}},$ $%
X_{i}\in \mathbb{R},$ $\mu \in \mathbb{R}.$ The likelihood of the sample is%
\begin{eqnarray*}
f(X_{1},...,X_{n};\mu ) &=&L(\mu )=\left( 2\pi \right) ^{-n/2}e^{-\frac{%
\sum_{i=1}^{n}\left( X_{i}-\mu \right) ^{2}}{2}} \\
\log L(\mu ) &=&-n/2\log \left( 2\pi \right) -\frac{1}{2}\tsum%
\nolimits_{i=1}^{n}\left( X_{i}-\mu \right) ^{2}
\end{eqnarray*}%
The MLE $\widehat{\mu }_{n}$ can be found as $\frac{\partial \log L(\mu )}{%
\partial \mu }=\sum_{i=1}^{n}\left( X_{i}-\mu \right) =0$ 
\begin{equation*}
\tsum\nolimits_{i=1}^{n}X_{i}=n\mu ,\text{ \ }\widehat{\mu }_{n}^{MLE}=\bar{X%
}_{n}
\end{equation*}
\end{example}

%TCIMACRO{\TeXButton{EndFrame}{\end{frame}}}%
%BeginExpansion
\end{frame}%
%EndExpansion

%TCIMACRO{\TeXButton{BeginFrame}{\begin{frame}}}%
%BeginExpansion
\begin{frame}%
%EndExpansion

\QTR{frametitle}{The Method of Maximum Likelihood}

\begin{example}[Normal with known variance]
Is $\widehat{\mu }_{n}^{MLE}$ unbiased and efficient? 
\begin{equation*}
\mathbb{E}\left( \widehat{\mu }_{n}^{MLE}\right) =\mathbb{E}\left( \frac{1}{n%
}\tsum\nolimits_{i=1}^{n}X_{i}\right) =\frac{1}{n}\tsum\nolimits_{i=1}^{n}%
\mathbb{E}\left( X_{i}\right) =\mu
\end{equation*}

To check if it is efficient, compute 
\begin{equation*}
\mathbb{V}\left( \widehat{\mu }_{n}^{MLE}\right) =\mathbb{V}\left( \frac{1}{n%
}\tsum\nolimits_{i=1}^{n}X_{i}\right) =\frac{1}{n^{2}}\tsum%
\nolimits_{i=1}^{n}\mathbb{V}\left( X_{i}\right) =\frac{1}{n}
\end{equation*}%
and compare that with the lower bound given by the $\left[ \mathcal{I}\left(
\mu \right) \right] ^{-1}=\left( -\mathbb{E}\left[ \frac{\partial ^{2}\log
L(\mu )}{\partial \mu ^{2}}\right] \right) ^{-1}=\left( -\mathbb{E}\left[ -n%
\right] \right) ^{-1}=\frac{1}{n},$ $\widehat{\mu }_{n}^{MLE}$ is efficient.
\end{example}

%TCIMACRO{\TeXButton{EndFrame}{\end{frame}}}%
%BeginExpansion
\end{frame}%
%EndExpansion

%TCIMACRO{\TeXButton{BeginFrame}{\begin{frame}}}%
%BeginExpansion
\begin{frame}%
%EndExpansion

\QTR{frametitle}{The Method of Maximum Likelihood}

\begin{example}[Poisson]
\emph{Let, }$X_{1},...,X_{n}$\emph{\ be i.i.d. Poisson }$\left( \lambda
\right) $\emph{\ r.v.s with mass }$f\left( X_{i},\lambda \right) =\frac{%
\lambda ^{X_{i}}}{X_{i}!}e^{-\lambda }$\emph{\ for }$X_{i}\in \left\{
0,1,2,...\right\} $\emph{. The likelihood function is given by }%
\begin{align*}
L\left( \lambda \right) & =\tprod\nolimits_{i=1}^{n}f\left( X_{i},\lambda
\right) =\tprod\nolimits_{i=1}^{n}\frac{\lambda ^{X_{i}}}{X_{i}!}e^{-\lambda
} \\
& =\frac{\lambda ^{X_{1}}...\lambda ^{X_{n}}}{\tprod\nolimits_{i=1}^{n}X_{i}!%
}e^{-n\lambda }=\frac{\lambda ^{X_{1}+...+X_{n}}}{\tprod%
\nolimits_{i=1}^{n}X_{i}!}.e^{-n\lambda }
\end{align*}
\end{example}

%TCIMACRO{\TeXButton{EndFrame}{\end{frame}}}%
%BeginExpansion
\end{frame}%
%EndExpansion

%TCIMACRO{\TeXButton{BeginFrame}{\begin{frame}}}%
%BeginExpansion
\begin{frame}%
%EndExpansion

\QTR{frametitle}{The Method of Maximum Likelihood}

\begin{example}[Poisson]
\emph{Taking logs:}%
\begin{equation*}
l\left( \lambda \right) =\log L\left( \lambda \right)
=\tsum\nolimits_{i=1}^{n}X_{i}\log \lambda -\log
\tprod\nolimits_{i=1}^{n}X_{i}!-n\lambda
\end{equation*}%
\begin{equation*}
\frac{\partial l}{\partial \lambda }=\frac{\sum_{i=1}^{n}X_{i}}{\lambda }%
-n=0\Rightarrow \lambda =\frac{1}{n}\tsum\nolimits_{i=1}^{n}X_{i}=\bar{X}%
_{n}.
\end{equation*}%
\emph{So the MLE of }$\lambda $ \emph{is }$\hat{\lambda}_{n}=\bar{X}_{n}.$
\end{example}

%TCIMACRO{\TeXButton{EndFrame}{\end{frame}}}%
%BeginExpansion
\end{frame}%
%EndExpansion

%TCIMACRO{\TeXButton{BeginFrame}{\begin{frame}}}%
%BeginExpansion
\begin{frame}%
%EndExpansion

\begin{example}[Poisson]
Now let's check if $\hat{\lambda}_{n}$ is unbiased, efficient or consistent.%
\begin{gather*}
\mathbb{E}[\hat{\lambda}_{n}]=\mathbb{E}[\frac{1}{n}\tsum%
\nolimits_{i=1}^{n}X_{i}]=\frac{1}{n}\tsum\nolimits_{i=1}^{n}\mathbb{E}%
[X_{i}] \\
\overset{\text{identically distributed}}{=}\frac{1}{n}\tsum%
\nolimits_{i=1}^{n}\mathbb{E}[X_{1}]=\mathbb{E}[X_{1}]. \\
\mathbb{E}[X]=e^{-\lambda }\tsum\nolimits_{j=0}^{\infty }j\frac{\lambda ^{j}%
}{j!}=\lambda e^{-\lambda }\tsum\nolimits_{j=1}^{\infty }\frac{\lambda ^{j-1}%
}{(j-1)!}=\lambda e^{-\lambda }e^{\lambda }=\lambda \\
\Rightarrow \mathbb{E}[\hat{\lambda}_{n}]=\lambda ,\text{ so }\hat{\lambda}%
_{n}\text{ is unbiased.}
\end{gather*}
\end{example}

%TCIMACRO{\TeXButton{EndFrame}{\end{frame}}}%
%BeginExpansion
\end{frame}%
%EndExpansion

%TCIMACRO{\TeXButton{BeginFrame}{\begin{frame}}}%
%BeginExpansion
\begin{frame}%
%EndExpansion

\begin{example}[Poisson]
Given that $\hat{\lambda}_{n}$ is unbiased, let's check if it is efficient.%
\begin{gather*}
var[\hat{\lambda}_{n}]=var[\frac{1}{n}\tsum\nolimits_{i=1}^{n}X_{i}]\overset{%
\text{independence}}{=}\frac{1}{n^{2}}\tsum\nolimits_{i=1}^{n}var[X_{i}] \\
\overset{\text{identically distributed}}{=}\frac{1}{n^{2}}%
\tsum\nolimits_{i=1}^{n}var[X_{1}]=\frac{1}{n}var[X_{1}]. \\
var[X]=\mathbb{E}[X^{2}]-\mathbb{E}^{2}[X] \\
\mathbb{E}[X(X-1)]=e^{-\lambda }\tsum\nolimits_{j=0}^{\infty }j(j-1)\frac{%
\lambda ^{j}}{j!}=\lambda ^{2}e^{-\lambda }\tsum\nolimits_{j=2}^{\infty }%
\frac{\lambda ^{j-2}}{(j-2)!} \\
\overset{i=j-2}{=}\lambda ^{2}e^{-\lambda }\tsum\nolimits_{i=0}^{\infty }%
\frac{\lambda ^{i}}{i!}=\lambda ^{2} \\
\Rightarrow \mathbb{E}[X^{2}]=\lambda ^{2}+\mathbb{E}[X] \\
var[X]=\lambda ^{2}+\lambda -\lambda ^{2}=\lambda
\end{gather*}
\end{example}

%TCIMACRO{\TeXButton{EndFrame}{\end{frame}}}%
%BeginExpansion
\end{frame}%
%EndExpansion

%TCIMACRO{\TeXButton{BeginFrame}{\begin{frame}}}%
%BeginExpansion
\begin{frame}%
%EndExpansion

\begin{example}[Poisson]
\begin{gather*}
\Rightarrow var[\hat{\lambda}_{n}]=\lambda /n \\
\mathcal{I}\left( \lambda \right) =-\mathbb{E}\left[ \frac{\partial ^{2}l}{%
\partial \lambda ^{2}}\right] =-\mathbb{E}\left[ -\frac{\sum_{i=1}^{n}X_{i}}{%
\lambda ^{2}}\right] =n/\lambda  \\
\Rightarrow var[\hat{\lambda}_{n}]=\mathcal{I}\left( \lambda \right) ^{-1},%
\text{ so }\hat{\lambda}_{n}\text{ is efficient.}
\end{gather*}%
Finally, to check if it consistent, 
\begin{equation*}
p\lim_{n\rightarrow \infty }\hat{\lambda}_{n}=p\lim_{n\rightarrow \infty }%
\frac{1}{n}\tsum\nolimits_{i=1}^{n}X_{i}\overset{\text{WLLN}}{=}\mathbb{E}%
[X_{1}]=\lambda 
\end{equation*}%
So, indeed, $\hat{\lambda}_{n}$ is consistent.
\end{example}

%TCIMACRO{\TeXButton{EndFrame}{\end{frame}}}%
%BeginExpansion
\end{frame}%
%EndExpansion

%TCIMACRO{\TeXButton{BeginFrame}{\begin{frame}}}%
%BeginExpansion
\begin{frame}%
%EndExpansion

\begin{example}[Rayleigh]
Suppose $X_{1},...,X_{n}$ is an i.i.d. sample of a Rayleigh Distribution
with density function: $f(X_{i};\zeta )=\frac{X_{i}}{\zeta ^{2}}\exp (-\frac{%
X_{i}^{2}}{2\zeta ^{2}}),$ $X_{i}\geq 0,$ $\zeta >0.$ The likelihood of the
sample is%
\begin{eqnarray*}
f(X_{1},...,X_{n};\zeta ) &=&L(\zeta )=\left[ \frac{1}{\zeta ^{2n}}%
\tprod\nolimits_{i=1}^{n}X_{i}\right] \exp \left( -\frac{1}{2\zeta ^{2}}%
\tsum\nolimits_{i=1}^{n}X_{i}^{2}\right) \\
\log L(\zeta ) &=&-2n\log \left( \zeta \right) +\tsum\nolimits_{i=1}^{n}\log
X_{i}-\frac{1}{2\zeta ^{2}}\tsum\nolimits_{i=1}^{n}X_{i}^{2}
\end{eqnarray*}%
The MLE $\widehat{\zeta }_{n}$ can be found as $\frac{\partial \log L(\zeta )%
}{\partial \zeta }=-\frac{2n}{\zeta }+\tsum\nolimits_{i=1}^{n}X_{i}^{2}\frac{%
1}{\zeta ^{3}}=0$ 
\begin{eqnarray*}
&\Longrightarrow &\frac{1}{\zeta }\left[ \frac{\tsum%
\nolimits_{i=1}^{n}X_{i}^{2}}{\zeta ^{2}}-2n\right] =0 \\
&\Longrightarrow &\widehat{\zeta }_{n}=\sqrt{\frac{\tsum%
\nolimits_{i=1}^{n}X_{i}^{2}}{2n}}\text{ since }\zeta >0
\end{eqnarray*}
\end{example}

%TCIMACRO{\TeXButton{EndFrame}{\end{frame}}}%
%BeginExpansion
\end{frame}%
%EndExpansion

%TCIMACRO{\TeXButton{BeginFrame}{\begin{frame}}}%
%BeginExpansion
\begin{frame}%
%EndExpansion

\begin{example}[Rayleigh]
Difficult to compute the bias of $\widehat{\zeta }_{n}$, but is it
consistent?%
\begin{gather*}
p\lim_{n\rightarrow \infty }\widehat{\zeta }_{n}=p\lim_{n\rightarrow \infty }%
\sqrt{\frac{\sum_{i=1}^{n}X_{i}^{2}}{2n}} \\
=\frac{1}{\sqrt{2}}\sqrt{p\lim_{n\rightarrow \infty }\left( \frac{%
\sum_{i=1}^{n}X_{i}^{2}}{n}\right) }\text{ by continuous mapping theorem} \\
\mathbb{E}[X^{2}]=\frac{1}{\zeta ^{2}}\int_{0}^{\infty }x^{3}\exp (-\frac{%
x^{2}}{2\zeta ^{2}})dx \\
=-\frac{1}{\zeta ^{2}}\int_{0}^{\infty }x^{3}\frac{2\zeta ^{2}}{2x}\frac{d}{%
dx}\left[ \exp (-\frac{x^{2}}{2\zeta ^{2}})\right] dx \\
=-\int_{0}^{\infty }x^{2}\frac{d}{dx}\left[ \exp (-\frac{x^{2}}{2\zeta ^{2}})%
\right] dx
\end{gather*}
\end{example}

%TCIMACRO{\TeXButton{EndFrame}{\end{frame}}}%
%BeginExpansion
\end{frame}%
%EndExpansion

%TCIMACRO{\TeXButton{BeginFrame}{\begin{frame}}}%
%BeginExpansion
\begin{frame}%
%EndExpansion

\begin{example}[Rayleigh]
\begin{gather*}
=-\left[ x^{2}\exp (-\frac{x^{2}}{2\zeta ^{2}})\right] _{0}^{\infty
}+2\int_{0}^{\infty }x\exp (-\frac{x^{2}}{2\zeta ^{2}})dx \\
=-\underset{\rightarrow 0}{\underbrace{\lim_{x\rightarrow \infty }\left[
x^{2}\exp (-\frac{x^{2}}{2\zeta ^{2}})\right] }}+\underset{=0}{\underbrace{%
\left[ x^{2}\exp (-\frac{x^{2}}{2\zeta ^{2}})\right] _{x=0}}} \\
+2\int_{0}^{\infty }\frac{\zeta ^{2}}{\zeta ^{2}}x\exp (-\frac{x^{2}}{2\zeta
^{2}})dx \\
=2\zeta ^{2}\underset{=1}{\underbrace{\int_{0}^{\infty }\frac{1}{\zeta ^{2}}%
x\exp (-\frac{x^{2}}{2\zeta ^{2}})dx}}=2\zeta ^{2} \\
\Rightarrow p\lim_{n\rightarrow \infty }\widehat{\zeta }_{n}=\frac{1}{\sqrt{2%
}}\sqrt{\mathbb{E}\left( X_{i}^{2}\right) }=\sqrt{\frac{2\zeta ^{2}}{2}}%
=\zeta ,\text{so, it is consistent.}
\end{gather*}
\end{example}

%TCIMACRO{\TeXButton{EndFrame}{\end{frame}}}%
%BeginExpansion
\end{frame}%
%EndExpansion

%TCIMACRO{\TeXButton{BeginFrame}{\begin{frame}}}%
%BeginExpansion
\begin{frame}%
%EndExpansion

\begin{example}[Rayleigh]
Let's compute the Fisher Information matrix. 
\begin{equation*}
\frac{\partial ^{2}\log L(\zeta )}{\partial \zeta ^{2}}=\frac{2n}{\zeta ^{2}}%
-3\sum_{i=1}^{n}X_{i}^{2}\frac{1}{\zeta ^{4}}
\end{equation*}%
\begin{equation*}
\Longrightarrow \mathcal{I}_{n}(\zeta )=-\mathbb{E}\left[ \frac{\partial
^{2}\log L(\zeta )}{\partial \zeta ^{2}}\right] =-\frac{2n}{\zeta ^{2}}+3%
\frac{1}{\zeta ^{4}}\sum_{i=1}^{n}\mathbb{E}[X_{i}^{2}]
\end{equation*}%
\QTR{Bbb}{E}$[X^{2}]=2\zeta ^{2},$ so%
\begin{equation*}
\mathcal{I}_{n}(\zeta )=-\frac{2n}{\zeta ^{2}}+3\frac{1}{\zeta ^{4}}2n\zeta
^{2}=\frac{4n}{\zeta ^{2}}
\end{equation*}
\end{example}

%TCIMACRO{\TeXButton{EndFrame}{\end{frame}}}%
%BeginExpansion
\end{frame}%
%EndExpansion

%TCIMACRO{\TeXButton{BeginFrame}{\begin{frame}}}%
%BeginExpansion
\begin{frame}%
%EndExpansion

\begin{example}[Rayleigh]
In the next two weeks, we shall prove an asymptotic normality result for any
MLE $\widehat{\theta }_{n}:$ $(\widehat{\theta }_{n}-\theta )\rightarrow _{d}%
\mathcal{N}(0,I_{n}(\theta )^{-1}).$

In this case that implies that $\widehat{\zeta }_{n}\rightarrow _{d}\mathcal{%
N}(\zeta ,\mathcal{I}_{n}(\zeta )^{-1}).$ Using this, we can also construct
95\% confidence intervals for the true $\zeta :$%
\begin{equation*}
\mathbb{P}\left( \widehat{\zeta }_{n}-1.96\sqrt{\frac{\zeta ^{2}}{4n}}\leq
\zeta \leq \widehat{\zeta }_{n}+1.96\sqrt{\frac{\zeta ^{2}}{4n}}\right)
\approx 95\%
\end{equation*}
\end{example}

%TCIMACRO{\TeXButton{EndFrame}{\end{frame}}}%
%BeginExpansion
\end{frame}%
%EndExpansion

%TCIMACRO{\TeXButton{BeginFrame}{\begin{frame}}}%
%BeginExpansion
\begin{frame}%
%EndExpansion

\QTR{frametitle}{The Method of Maximum Likelihood}

Let $X_{1},...,X_{n}$\ be i.i.d. $\mathcal{N}\left( \mu ,\sigma ^{2}\right) $%
\ random variables. Here, we have a vector of unknown parameters: $\theta
=\left( \mu ,\sigma ^{2}\right) ^{\prime }$.%
\begin{align*}
L\left( \mu ,\sigma ^{2}\right) & =\prod\limits_{i=1}^{n}f\left( X_{i};\mu
,\sigma ^{2}\right) =\prod\limits_{i=1}^{n}\left( 2\pi \sigma ^{2}\right)
^{-1/2}\exp \left\{ -\frac{1}{2\sigma ^{2}}\left( X_{i}-\mu \right)
^{2}\right\}  \\
& =\left( 2\pi \sigma ^{2}\right) ^{-n/2}\exp \left\{ -\frac{1}{2\sigma ^{2}}%
\sum_{i=1}^{n}\left( X_{i}-\mu \right) ^{2}\right\} ,
\end{align*}%
so the log-likelihood is given by%
\begin{equation*}
l\left( \mu ,\sigma ^{2}\right) =-\frac{n}{2}\log 2\pi -\frac{n}{2}\log
\sigma ^{2}-\frac{1}{2\sigma ^{2}}\sum_{i=1}^{n}\left( X_{i}-\mu \right)
^{2}.
\end{equation*}

%TCIMACRO{\TeXButton{EndFrame}{\end{frame}}}%
%BeginExpansion
\end{frame}%
%EndExpansion

%TCIMACRO{\TeXButton{BeginFrame}{\begin{frame}}}%
%BeginExpansion
\begin{frame}%
%EndExpansion

\QTR{frametitle}{The Method of Maximum Likelihood}

\begin{example}
The MLE for $\left( \mu ,\sigma ^{2}\right) ^{\prime }$\ will be found by
jointly maximising $l\left( \mu ,\sigma ^{2}\right) $\ for $\mu $\ and $%
\sigma ^{2}$, i.e. it will be the solution of the two-equation system $\frac{%
\partial l}{\partial \mu }=0$\ and $\frac{\partial l}{\partial \sigma ^{2}}%
=0 $. Since $\sigma ^{2}>0$\ 
\begin{equation*}
\frac{\partial l}{\partial \mu }=\frac{1}{\sigma ^{2}}\sum_{i=1}^{n}\left(
X_{i}-\mu \right) =\frac{1}{\sigma ^{2}}\left( n\bar{X}_{n}-n\mu \right)
=0\Longrightarrow \mu =\bar{X}_{n}.
\end{equation*}

Differentiating the log-likelihood w.r.t. $\sigma ^{2}:$%
\begin{align*}
\frac{\partial l}{\partial \sigma ^{2}}& =-\frac{n}{2}\frac{1}{\sigma ^{2}}+%
\frac{1}{2\left( \sigma ^{2}\right) ^{2}}\sum_{i=1}^{n}\left( X_{i}-\mu
\right) ^{2}=0 \\
& \Longrightarrow -n+\frac{1}{\sigma ^{2}}\sum_{i=1}^{n}\left( X_{i}-\mu
\right) ^{2}=0
\end{align*}
\end{example}

%TCIMACRO{\TeXButton{EndFrame}{\end{frame}}}%
%BeginExpansion
\end{frame}%
%EndExpansion

%TCIMACRO{\TeXButton{BeginFrame}{\begin{frame}}}%
%BeginExpansion
\begin{frame}%
%EndExpansion

\QTR{frametitle}{The Method of Maximum Likelihood}

\begin{example}
$\Longrightarrow \sigma ^{2}=\frac{1}{n}\sum_{i=1}^{n}\left( X_{i}-\mu
\right) ^{2}.$

Requiring the 2 F.O.C.s to hold simultaneously, we obtain the MLE $\left( 
\hat{\mu}_{n},\hat{\sigma}_{n}^{2}\right) $\ of $\left( \mu ,\sigma
^{2}\right) $ to be $\hat{\mu}_{n}=\bar{X}_{n},\;\;\;\;\hat{\sigma}_{n}^{2}=%
\frac{1}{n}\sum_{i=1}^{n}\left( X_{i}-\bar{X}_{n}\right) ^{2}.$
\end{example}

\begin{remark}

$\hat{\sigma}_{n}^{2}$ is biased. Thus, MLEs are not unbiased in general.
Despite it being biased, the MLE estimator has smaller MSE than the unbiased
estimator $s_{n}^{2}$.

\end{remark}

%TCIMACRO{\TeXButton{EndFrame}{\end{frame}}}%
%BeginExpansion
\end{frame}%
%EndExpansion

%TCIMACRO{\TeXButton{BeginFrame}{\begin{frame}}}%
%BeginExpansion
\begin{frame}%
%EndExpansion

\QTR{frametitle}{The Method of Maximum Likelihood}

\begin{itemize}
\item When $X_{1},..,X_{n}$\ are not independent, $L\left( \theta \right)
\neq \prod\limits_{i=1}^{n}f\left( X_{i};\theta \right) $. A way to proceed
with setting up the likelihood function is to use the definition of
conditional density:%
\begin{equation*}
f\left( X_{2}|X_{1}\right) =\frac{f\left( X_{2},X_{1}\right) }{f\left(
X_{1}\right) }\Longrightarrow f\left( X_{2},X_{1}\right) =f\left(
X_{2}|X_{1}\right) f\left( X_{1}\right) .
\end{equation*}%
\begin{align*}
f\left( X_{3},X_{2},X_{1}\right) & =f\left( X_{3}|X_{2},X_{1}\right) f\left(
X_{2},X_{1}\right)  \\
& =f\left( X_{3}|X_{2},X_{1}\right) f\left( X_{2}|X_{1}\right) f\left(
X_{1}\right) 
\end{align*}

\item Proceeding like this we obtain the joint density of $\left(
X_{n},...,X_{1}\right) $:%
\begin{eqnarray*}
f\left( X_{n},X_{n-1},...,X_{1}\right)  &=&f\left(
X_{n}|X_{n-1},...,X_{1}\right) f\left( X_{n-1}|X_{n-2},...,X_{1}\right) ...
\\
&&\times ...f\left( X_{2}|X_{1}\right) f\left( X_{1}\right) .
\end{eqnarray*}
\end{itemize}

%TCIMACRO{\TeXButton{EndFrame}{\end{frame}}}%
%BeginExpansion
\end{frame}%
%EndExpansion

%TCIMACRO{\TeXButton{BeginFrame}{\begin{frame}}}%
%BeginExpansion
\begin{frame}%
%EndExpansion

\QTR{frametitle}{The Method of Maximum Likelihood}

The above factorisation yields the following expressions for the likelihood
and log-likelihood functions: 
\begin{equation*}
L\left( \theta \right) =f\left( X_{1};\theta \right)
\prod\limits_{i=2}^{n}f\left( X_{i}|X_{i-1},..,X_{1};\theta \right) ,
\end{equation*}%
\begin{equation*}
l\left( \theta \right) =\log f\left( X_{1};\theta \right)
+\sum_{i=2}^{n}\log f\left( X_{i}|X_{i-1},..,X_{1};\theta \right) .
\end{equation*}

%TCIMACRO{\TeXButton{EndFrame}{\end{frame}}}%
%BeginExpansion
\end{frame}%
%EndExpansion

\section{Reading}

%TCIMACRO{\TeXButton{BeginFrame}{\begin{frame}}}%
%BeginExpansion
\begin{frame}%
%EndExpansion

\QTR{frametitle}{Reading}

\begin{itemize}
\item Relevant chapters from the textbook: 5
\end{itemize}

%TCIMACRO{\TeXButton{EndFrame}{\end{frame}}}%
%BeginExpansion
\end{frame}%
%EndExpansion

\end{document}
